% Solutions to the Chapter 9 exercises that are not code, from lectures/L09/appendix/c_exercises.md.
% Exercises 9.6 and 9.7 are Code and are not answered; see chapter.tex for why.

\section{Chapter 9: Sleeping, Waiting, and Time}
\label{sol:sec:c9}
Answers to the exercises in \secref{c9:sec:exercises}. Exercises~\ref{c9:ex:6} and~\ref{c9:ex:7}
are Code, and are checked on the target as their Check yourself lines say.

% ------------------------------------------------------------------------------------------
\solution[fillaccent2]{c9:ex:1}{Sleeping and spinning}{Recall}

\xpart{a} The sleeping task is in \code{TASK_INTERRUPTIBLE}, \code{S} in \code{ps}, or
\code{TASK_UNINTERRUPTIBLE}, \code{D}; it is not on the run queue, and it costs no CPU time at all,
only its memory and the wake-up when it comes. The spinning task is in \code{TASK_RUNNING},
\code{R}; it is on the run queue, and it costs a whole CPU while it has one, every time slice taken
from a task with work to do (\secref{c9:app:a1}).

\xpart{b} A signal, even \code{SIGKILL}, is acted on only when the task is on its way back to
userspace, and a task in \code{TASK_UNINTERRUPTIBLE} is not woken by signals at all, only by an
explicit wake. So the signal waits, pending, until whatever the task is waiting for finishes. That
is usually I/O: a disk or network-filesystem request, a page being read in, a device operation. A
\code{D} in \code{ps} says something the task depends on is not completing, a stuck device, a hung
driver, a server that has gone away; the process is the symptom, and the cause is below it. On
Linux such a task also counts in the load average.

\xpart{c} \code{TASK_INTERRUPTIBLE}. A user can always interrupt or kill the reader, and your
driver answers with \code{-ERESTARTSYS}. Put a reader in \code{TASK_UNINTERRUPTIBLE} to wait for a
device that never produces data, and you have made the unkillable \code{D} process of part
\textbf{b)} yourself. The uninterruptible state is for short waits that cannot be abandoned half
way. (For a wait that only a fatal signal should end, there is a third, \code{TASK_KILLABLE}, which
\code{wait_event_killable} uses.)

\xpart{d} The task's state must be set to something other than \code{TASK_RUNNING} before the call.
With the state still \code{TASK_RUNNING}, \code{schedule()} leaves the task on the run queue and
may pick it again at once: a yield. And for the sleep ever to end, the task must be somewhere a
wake-up can find it, on the wait queue the waker uses.

% ------------------------------------------------------------------------------------------
\solution[fillaccent2]{c9:ex:2}{The lost wakeup}{Recall}

\xpart{a}
\begin{enumerate}
  \item The reader calls \code{fifo_empty()}, which finds the FIFO empty and returns true.
  \item \textbf{The interrupt lands here}, after the test and before \code{set_current_state}. The
    handler pushes a sample and calls \code{wake_up_interruptible}. The reader is on the queue but
    still \code{TASK_RUNNING}, so there is nothing to wake, and the wake does nothing.
  \item The reader sets \code{TASK_INTERRUPTIBLE} and calls \code{schedule()}, and leaves the run
    queue with a sample waiting for it.
  \item Nothing wakes it until the next \code{wake_up}: on a device that interrupts again, a period
    later; if that was the last event, never.
\end{enumerate}

\xpart{b} Two instructions: the load inside \code{fifo_empty()} that finds the FIFO empty, and the
store in \code{set_current_state}. That is nanoseconds, and how soon it is hit is the chance of one
attempt times the attempts per second. An interrupt at an unrelated moment lands in a window $w$
wide, in a period of 1 ms, with probability $w / 1$ ms, of the order of a few in a million for a few
nanoseconds; and a reader that blocks for each sample offers the window up to a thousand times a
second. That is a few chances in a thousand every second: on paper, minutes of continuous running.
On a real machine it depends on how long the window really lasts and on when the reader's sleeps
fall relative to the interrupts, and it is found by running it. Either way it is found in a test
run, not once in a product's lifetime.

\xpart{c} \code{long __int = prepare_to_wait_event(&wq_head, &__wq_entry, state);}. It puts the task
on the queue and sets its state, under the queue's lock, before the condition is tested; a
\code{wake_up} that comes after it finds a task to wake (\secref{c9:app:a3}).

\xpart{d} \textbf{The lost wakeup}: the test has to happen after the task is queued, so the macro
must be able to evaluate it there, inside the loop; a variable's value passed in would be read once,
at the call, before queueing, which is exactly the bug. \textbf{After a wake}: a wake does not mean
the condition is true, because another reader may have taken the data or the wake may be spurious,
so the macro evaluates the expression again every time round the loop.

\xpart{e} It must be \textbf{cheap}, and \textbf{free of side effects}, because it is evaluated once
before the task is queued and again after every wake, as many times as that takes. It is also
evaluated with the task already in \code{TASK_INTERRUPTIBLE}, so it must not sleep: a sleeping call
there resets the state, and \code{CONFIG_DEBUG_ATOMIC_SLEEP} warns \code{do not call blocking ops
when !TASK_RUNNING}.

\xpart{f} It measures the wrong thing. What decides how often the bug fires is not the width of the
window but the chance of one attempt times the number of attempts: the window is offered every time
a reader sleeps, a thousand times a second, on every unit shipped, for years, and a chance that is
rare per attempt is certain in aggregate. Rare is also the worse outcome: a bug that fires once a
week passes every test and cannot be reproduced when it is reported. And a window you can point at
is a bug whatever its width, when the fix, queueing before testing, costs nothing.

% ------------------------------------------------------------------------------------------
\solution[fillaccent3]{c9:ex:3}{Jiffies}{Hand calculation}

\xpart{a} $1{,}000{,}000 / 250 = \mathbf{4{,}000}$ us, 4 ms.

\xpart{b} With \code{HZ} dividing 1,000 exactly, \code{msecs_to_jiffies(m)} is $(m + 3) / 4$ in
integer arithmetic: $(1 + 3)/4 = 1$, $(4 + 3)/4 = 1$ and $(10 + 3)/4 = 3$. It rounds \textbf{up},
and that matters because a timeout must never be shorter than it was asked to be. Rounded down,
\code{msecs_to_jiffies(1)} would be 0, and \code{msleep(1)}, which sleeps while its count of
jiffies is not zero, would not sleep at all; and 10 ms would become 2 jiffies, 8 ms.

\xpart{c} One jiffy more than \code{msecs_to_jiffies(n)}: \code{msleep(1)} is $(1 + 1) \times 4 =
\mathbf{8}$ ms, \code{msleep(4)} is also $(1 + 1) \times 4 = \mathbf{8}$ ms, and \code{msleep(10)} is
$(3 + 1) \times 4 = \mathbf{16}$ ms. The course measured 8,000, 8,000 and 15,998 us
(\secref{c9:app:b4}). The extra jiffy is not in \code{msleep}, which passes
\code{msecs_to_jiffies(n)} to \code{schedule_timeout} unchanged; it is the timer wheel's, which
rounds every expiry up by a jiffy so that a timer armed part of the way through a tick never fires
early (\code{calc_index} in \file{kernel/time/timer.c}). Each call in a loop starts just after the
tick that woke the one before, so it sleeps $n + 1$ whole jiffies.

\xpart{d} \code{msleep(1)} and \code{msleep(4)}, 8 ms each. The loop runs eight times slower than
its author believes: 1,000 iterations take 8 s, not 1 s. Asking for less than a jiffy buys nothing.

\xpart{e} At \code{HZ=1000} a jiffy is 1 ms and \code{msecs_to_jiffies(n)} is $n$: \code{msleep(1)}
is $1 + 1 = 2$ ms, \code{msleep(4)} is $4 + 1 = 5$ ms, and \code{msleep(10)} is $10 + 1 = 11$ ms.

\xpart{f} \code{>} breaks when the counter wraps between computing the deadline and comparing
against it. Computed 50 ticks before the wrap, \code{jiffies + 125} wraps round to 75, and
\code{jiffies > timeout} is true at once: the timeout expires immediately. Computed so that
\code{jiffies} wraps first, the comparison is false until the counter comes all the way round, and
the timeout never expires. \code{time_after} compares the signed difference, which is right as long
as the two are less than half the range apart. The wrap comes every $2^{32} / 250 = 17{,}179{,}869$
s, \textbf{198.8 days}; but the first one comes \textbf{five minutes after boot}, because the kernel
starts the count at \code{INITIAL_JIFFIES}, $-300 \times$ \code{HZ}, on purpose, so that code with
this bug fails in its first five minutes rather than after half a year. (On the course's 64-bit
target \code{jiffies} itself never wraps; its low 32 bits still do, five minutes after boot, which
catches code that keeps it in a \code{u32}.)

% ------------------------------------------------------------------------------------------
\solution[fillaccent3]{c9:ex:4}{Choosing a delay}{Hand calculation}

\xpart{a} \textbf{\code{udelay(2)}.} Inside a spinlock nothing may sleep, and 2 us of busy-waiting
costs less than any sleep would to arrange.

\xpart{b} \textbf{\code{usleep_range}} between reads of the bit, for example
\code{usleep_range(100, 200)}, with an overall timeout; \code{readl_poll_timeout} is that loop
already written. \code{probe} may sleep, \code{udelay} would burn the whole 500 us, and
\code{msleep(1)} would cost 8 ms a poll on this kernel.

\xpart{c} \textbf{\code{msleep(100)}}, in process context: \code{msecs_to_jiffies(100)} is 25, so
the call takes 26 jiffies, 104 ms, and a quantum of 4 ms is noise at that length. If the retry must
not block its caller, a timer.

\xpart{d} \textbf{None of them} (this is the trick). A hard interrupt handler may not sleep, and a
busy-wait of 5 ms, \code{mdelay(5)}, is legal but spends 5 ms with interrupts disabled on that CPU.
Do not wait in the handler: acknowledge, and do the waiting in the threaded half, or arm a timer to
do the next step 5 ms later.

\xpart{e} \textbf{A timer}, a \code{timer_list} armed 2 s ahead and cancelled when the device
answers: a timeout that is usually cancelled and whose precision is irrelevant is the case the
wheel exists for (\secref{c9:app:b5}).

\xpart{f} \textbf{\code{usleep_range(50, 50)}}: process context may sleep, and with no slack an
hrtimer is the most accurate way to. The alternative is \code{udelay(50)}, and the course's
measurements are its argument: \code{usleep_range(100, 100)} took 307 us and \code{udelay(100)} 112
(\secref{c9:app:b4}), so on this target a sleep carries about 200 us of wake-up overhead. If the
requirement is 50 us and not 250, \code{udelay} is the one that delivers it, at the price of 50 us
of a CPU every time.

\xpart{g} \textbf{d)}: of the four ways to wait it offers, none is right in a hard interrupt
handler, and the right answer is not to wait there.

% ------------------------------------------------------------------------------------------
\solution[fillaccent4]{c9:ex:5}{The four cases of \code{read}}{Design}

\xpart{a} \textbf{160.} The count is in bytes and a sample is 32 bits, so 100 samples is 400 bytes,
and \code{read} returns the 40 samples there are, 160 bytes. The program sees a short read at once,
and calls \code{read} again if it wants more. (Had the 100 been bytes, 25 whole samples: 100.)

\xpart{b} \textbf{Nothing yet}: the driver sleeps in \code{wait_event_interruptible} until the
handler adds a sample and wakes it, then returns the whole samples available, at least 4 bytes. The
program sees \code{read} block for up to a period, about a millisecond on a running \code{qa-dev},
and then return a positive count.

\xpart{c} \textbf{\code{-EAGAIN}}, at once. The program sees $-1$ with \code{errno} set to
\code{EAGAIN}, ``Resource temporarily unavailable''.

\xpart{d} \textbf{\code{-ERESTARTSYS}}: the wait returns it because a signal is pending, and the
driver passes it on. The program never sees it. Ctrl-C sends \code{SIGINT}, whose default action
ends the process, so \code{read} never returns and the program dies, which is what Ctrl-C is for. A
program with a handler installed with \code{SA_RESTART} runs the handler and has its \code{read}
restarted without noticing; one whose handler has no \code{SA_RESTART} sees $-1$ with \code{errno}
set to \code{EINTR}.

\xpart{e} The driver now has something that can wake a reader, the interrupt handler of
Chapter~\ref{c8:ch}, and a way to sleep without losing the wake. Chapter~\ref{c5:ch} had neither:
it could not wait correctly, so the only honest answers were the data or ``nothing now'', and 0
would have said ``nothing ever''.

\xpart{f} \textbf{\code{-EINTR}} takes the decision away from the kernel. The program sees
\code{EINTR} whenever a handled signal arrives, even one installed with \code{SA_RESTART}, which
asked for the call to be restarted; a program that does not retry on \code{EINTR} then fails for
no reason it can see, whenever a \code{SIGCHLD} or a timer signal happens by. \textbf{Ignoring the
return value} depends on what comes next. If the code goes on to take data from a FIFO that is
still empty, it returns 0, end of file, and the program stops as if the device had finished. If it
loops back to wait, the signal is still pending, so every wait returns at once and the
\code{read} spins in the kernel at full speed without ever returning, and a signal is delivered
only on the way back: a reader that cannot be killed (\secref{c9:app:a:-erestartsys}).

\xpart{g} End of file: no more data, ever. \code{cat /dev/qa_wait} exits at once with status 0
having printed nothing, as if the device had finished, and so does any loop that reads until
\code{read} returns 0.

% ------------------------------------------------------------------------------------------
\solution[fillaccent4]{c9:ex:8}{A timer instead of an interrupt}{Design}

\xpart{a} Load the module with \code{ENABLE} set and \code{IRQ_EN} clear, and arm a
\code{timer_list}. Its callback reads the device's FIFO, puts the samples into the driver's FIFO
under the spinlock, wakes the readers, and re-arms itself. The callback runs in softirq context
(\secref{c9:app:b5}): it may touch registers, take a spinlock and wake a queue, and may not sleep,
take a mutex, allocate with \code{GFP_KERNEL} or call \code{copy_to_user}. What moves: the lock's
process-context side must now disable softirqs as well, \code{spin_lock_bh} or the \code{irqsave}
it already has, and anything that sleeps stays in \code{read}. Two details. 4 ms is one jiffy at
\code{HZ=250}, the finest a \code{timer_list} does, and it must be re-armed from its previous
expiry, \code{mod_timer(&t, t.expires + 1)}: re-armed from \code{jiffies + 1} it polls every
8 ms, for the same reason \code{msleep(1)} takes two jiffies. And on unload it re-arms itself, so
it is stopped with \code{timer_shutdown_sync}, which also refuses any later re-arm.

\xpart{b} The resolution is no longer the jiffy: any period, 4 ms or 1 ms, exactly, and periodic
from its previous expiry with \code{hrtimer_forward_now} and \code{HRTIMER_RESTART}. The callback
runs in hard interrupt context, which is the default on this kernel, so it is more restricted than
the \code{timer_list}'s: the lock shared with \code{read} must be taken \code{irqsave}, and its time
is interrupt time again. (A \code{_SOFT} mode moves it to softirq, and under \code{PREEMPT_RT} any
hrtimer not marked \code{_HARD} is moved there.)

\xpart{c} With the \code{timer_list} the sample waits in the device's FIFO until the next poll, up
to 4 ms, more if the timer is late. With the hrtimer, up to its period, which can be as short as the
wake-ups you are willing to pay for. With the interrupt it is taken when it arrives. Polling adds,
on average, half its period to the latency of every sample.

\xpart{d} Sixteen samples at one a millisecond fill the FIFO in 16 ms, and the seventeenth, 16 ms
after the first undrained one, is dropped: two polls must never be more than \textbf{16 ms} apart,
and the nominal period has to leave room for how late the timer can be. With a FIFO of 4 it is
\textbf{4 ms}, exactly one jiffy, which a \code{timer_list} cannot promise, since it may fire a
jiffy late; that needs an hrtimer, or the interrupt.

\xpart{e} Poll when you know when the data will be there, or when it is always there: a steady
stream in which every poll finds work, which is why network drivers switch to polling under load.
Poll too when the device cannot interrupt. Use the interrupt when events are sporadic, or when
latency matters: a poll that usually finds nothing is pure cost, and its period is added latency. A
timer is what you use when nothing tells you (\secref{c9:app:b:choosing}).

% ------------------------------------------------------------------------------------------
\solution[fillaccent]{c9:ex:9}{A millisecond that is not a millisecond}{Cross-check}

\xpart{a} One jiffy is 4 ms, and each \code{msleep} is one jiffy more than \code{msecs_to_jiffies}
gives:

\begin{narrowtable}{@{}llr@{}}
  \thead{Call} & \thead{Jiffies} & \thead{Predicted} \\
  \midrule
  \code{msleep(1)} & $1 + 1 = 2$ & 8,000 us \\
  \code{msleep(4)} & $1 + 1 = 2$ & 8,000 us \\
  \code{msleep(10)} & $3 + 1 = 4$ & 16,000 us \\
  \code{usleep_range(1000, 1000)} & not quantised & a little over 1,000 us \\
\end{narrowtable}

\noindent \code{usleep_range} uses an hrtimer, so it costs the request plus the time to take the
timer interrupt and switch the task back in; the model gives no figure for that.

\xpart{c}
\begin{itemize}
  \item \code{msleep(1)}: 8,000 us for 1,000 asked, \textbf{eight times}. The predictions of part
    \textbf{a)} expect it exactly; a prediction from ``at least one jiffy'' alone says 4 ms, out by
    two. Followed as written, the prediction furthest from its measurement is
    \code{usleep_range}'s: 1,000 us against 1,443, a factor of 1.44 (\secref{c9:app:b4}).
  \item \code{msleep(1)} and \code{msleep(4)}, 8,000 us each. Not a bug in the measurement or in
    \code{msleep}, but in the expectation: \code{msleep} promises at least what it was asked, in
    jiffies, and 1 ms and 4 ms are both one jiffy.
  \item \code{msleep} passes \code{msecs_to_jiffies(n)} to \code{schedule_timeout} with nothing
    added. The sentence is ``at least \code{@timeout} jiffies are guaranteed to pass before the
    routine returns'', from its comment in \file{kernel/time/timer.c}, for the uninterruptible
    sleep \code{msleep} uses. The call arrives part of the way through a jiffy, and counting that
    partial jiffy as one would return early, so the timer must fire at the $(n + 1)$th tick after
    the call; the timer wheel does it by rounding every expiry up by one jiffy. In a loop each call
    starts just after a tick, and so waits $n + 1$ whole jiffies.
  \item A high-resolution timer, programmed on the hardware timer rather than on the jiffy wheel.
    What it gave up is the range: \code{usleep_range(1000, 1000)} leaves the kernel no slack, so it
    must program a wake-up for this sleep alone rather than fold it into one already due, which is
    the saving the range exists for.
\end{itemize}

\xpart{d} \code{msecs_to_jiffies(5)} is $(5 + 3)/4 = 2$, so \code{msleep(5)} takes $2 + 1 = 3$
jiffies, \textbf{12,000 us}. That is the test that separates the accounts: ``one jiffy more than
\code{msecs_to_jiffies}'' predicts 12 ms, where ``twice the rounded-up request'' would predict 16.
The course does not publish this measurement; the prediction is 12,000 us.

\xpart{e} At \code{HZ=1000}: \code{msleep(1)} 2 ms, \code{msleep(4)} 5 ms, \code{msleep(10)} 11 ms,
and \code{usleep_range(1000, 1000)} a little over 1 ms, as before. The three \code{msleep}s change
and \code{usleep_range} does not. \code{msleep} is quantised: it costs its request rounded up to
jiffies, plus one, and the jiffy is \code{HZ}'s. \code{usleep_range} is not: it costs its request
plus a wake-up overhead that has nothing to do with \code{HZ}.

\xpart{f} On this kernel $1{,}000 \times 8$ ms, \textbf{8 s}; at \code{HZ=1000}, $1{,}000 \times 2$
ms, \textbf{2 s}. Written as \code{msleep(1000);} it takes one second on both, to within the timer
wheel's rounding. At that length the wheel places the timer in a coarser level and rounds its expiry
up to that level's step: 8 jiffies at \code{HZ=250}, so between 1,000 and 1,032 ms, and 64 jiffies
at \code{HZ=1000}, so between 1,000 and 1,064 ms. Either is a second, for a comment.

\xpart{g} Put what was measured and how: average durations over 200 repetitions on an idle QEMU
guest, kernel 6.12.30 with \code{PREEMPT} and \code{HZ=250}, and the structure that is not a
measurement, that \code{msleep(n)} waits at least $n$ ms and in a loop costs
\code{msecs_to_jiffies(n)} plus one jiffies. Refuse to claim any maximum or worst case, and any
figure under load or on other hardware: an average over an idle machine says nothing about the
tail, and these numbers contain the emulator. Measure instead the whole response, from the event to
the driver's action, as a distribution over a long run under realistic load, on the product's own
hardware and configuration, and report its high percentiles and the worst value seen, with a trace
of what caused the outliers; and even then a measured maximum is an observation, not a bound.
Chapter~\ref{c12:ch} is that measurement.
